\section{Functional calculus on complementary projection ranges}
\label{sec:projection_calculus}

\begin{theorem}[Functional calculus of an affine projection]
    \label{thm:area_law_affine_projection_calculus}
    \lean{Matrix.cfc_affine_of_isSelfAdjoint_isIdempotentElem}
    Let $P$ be an orthogonal projection on a finite-dimensional complex
    Hilbert space. For arbitrary $b,d\in\mathbb R$ and an arbitrary
    function $f:\mathbb R\to\mathbb R$,
    \begin{align}
        f\bigl(bI+(d-b)P\bigr) & =f(b)(I-P)+f(d)P.
        \label{eq:affine_projection_calculus}
    \end{align}
    Here functional calculus uses the restriction of $f$ to the finite
    spectrum. Neither continuity of $f$ nor distinctness of $b,d$ is required.
    Either projection range, or the entire space, may have dimension zero.
    This is the complementary-range calculation used in
    \cite[Section~7, lines~603--640]{OpenAI2026AreaLaw}.
\end{theorem}

\begin{proof}
    \uses{lem:orthogonal_resolution_cfc}
    The projections $I-P,P$ are orthogonal and sum to $I$. Apply
    Lemma~\ref{lem:orthogonal_resolution_cfc} to this resolution with scalar
    values $b,d$. This argument permits zero projections and repeated scalar
    values. In particular, if $b=d$, the formula reduces to $f(b)I$; if
    $P=0$ or $P=I$, it reduces to $f(b)I$ or $f(d)I$, respectively.
\end{proof}
